% GENERATED FILE. Edit canonical lessons, then run npm run generate:books.
% canonical-source-hash: fnv1a64:87a62e5f19c96192
% canonical-lessons: FR-C515-trac, FR-C515-replique, FR-C515-interpreter, FR-C515-comedien, FR-C515-applaudissements

\chapter{Sur scène}
\label{ch:fr-sur-scene}
\hlchaptermodality{515}

\begin{chapteropening}
\textbf{By the end of this chapter:} \emph{I can talk about acting on stage.}

Complete the last lesson of chapter 515: I can talk about acting on stage.
\end{chapteropening}

\section[le trac]{le trac — stage fright, nerves}

\label{lesson:FR-C515-trac}



\begin{quote}
Before the new one: say the French for costume, then the French for interval.
\end{quote}

\subsection*{You'll want to know: le trac}
\textbf{le trac} — \textquotedblleft{}stage fright, nerves\textquotedblright{}.

\begin{quote}
Avant chaque concert, elle a le trac.
\end{quote}

\begin{quote}
\emph{Before every concert she gets stage fright.}
\end{quote}

\emph{Avoir le trac} — to have nerves, also before an exam.

\begin{cousinweb}
The origin of \emph{le trac} is not settled; it appeared in French in the nineteenth century.
\end{cousinweb}

\begin{grammarlens}[title={What you've built}]
One more word for telling, explaining and describing at B1.
\end{grammarlens}

\subsection*{Your turn}
\begin{itemize}
\raggedright
  \item \emph{Say it:} \emph{le trac}
  \item \emph{Say it:} \emph{le trac}, once more
  \item \emph{Say it:} a sentence of your own with \emph{le trac}
  \item \emph{Recall:} say the French for costume, then the French for interval, then say \emph{le trac} again
\end{itemize}

\begin{tcolorbox}[breakable,colback=teal!4,colframe=teal!35!black,title={Before you move on}]
What does le trac mean? (\textquotedblleft{}Stage fright, nerves\textquotedblright{}.) Say it once more.
\end{tcolorbox}
\section[la réplique]{la réplique — line (in a play); reply}

\label{lesson:FR-C515-replique}



\begin{quote}
Before the new one: say the French for interval, then the French for stage fright.
\end{quote}

\subsection*{You'll want to know: la réplique}
\textbf{la réplique} — \textquotedblleft{}line (in a play); reply\textquotedblright{}.

\begin{quote}
L'acteur a oublié sa réplique.
\end{quote}

\begin{quote}
\emph{The actor forgot his line.}
\end{quote}

Also a copy: \emph{une réplique de la statue}.

\begin{cousinweb}
\emph{la réplique} is built on \emph{répliquer} \textquotedblleft{}to reply\textquotedblright{}, from Latin \emph{replicāre} \textquotedblleft{}to fold back, to reply\textquotedblright{}, built on \emph{plicāre} \textquotedblleft{}to fold\textquotedblright{} (Lewis \& Short).
\end{cousinweb}

\begin{grammarlens}[title={What you've built}]
One more word for telling, explaining and describing at B1.
\end{grammarlens}

\subsection*{Your turn}
\begin{itemize}
\raggedright
  \item \emph{Say it:} \emph{la réplique}
  \item \emph{Say it:} \emph{la réplique}, once more
  \item \emph{Say it:} a sentence of your own with \emph{la réplique}
  \item \emph{Recall:} say the French for interval, then the French for stage fright, then say \emph{la réplique} again
\end{itemize}

\begin{tcolorbox}[breakable,colback=teal!4,colframe=teal!35!black,title={Before you move on}]
What does la réplique mean? (\textquotedblleft{}Line (in a play); reply\textquotedblright{}.) Say it once more.
\end{tcolorbox}
\section[interpréter]{interpréter — to perform, to play (a role); to interpret}

\label{lesson:FR-C515-interpreter}



\begin{quote}
Before the new one: say the French for stage fright, then the French for line.
\end{quote}

\subsection*{You'll want to know: interpréter}
\textbf{interpréter} — \textquotedblleft{}to perform, to play (a role); to interpret\textquotedblright{}.

\begin{quote}
Elle interprète le rôle de la mère.
\end{quote}

\begin{quote}
\emph{She plays the role of the mother.}
\end{quote}

Also to interpret a dream or a text. An interpreter between languages is \emph{un interprète}.

\begin{cousinweb}
\emph{interpréter} is a learned borrowing from Latin \emph{interpretārī} \textquotedblleft{}to explain, to translate\textquotedblright{} (EtymDB), from \emph{interpres} \textquotedblleft{}a go-between, an interpreter\textquotedblright{} (Lewis \& Short).
\end{cousinweb}

\begin{grammarlens}[title={What you've built}]
One more word for telling, explaining and describing at B1.
\end{grammarlens}

\subsection*{Your turn}
\begin{itemize}
\raggedright
  \item \emph{Say it:} \emph{interpréter}
  \item \emph{Say it:} \emph{interpréter}, once more
  \item \emph{Say it:} a sentence of your own with \emph{interpréter}
  \item \emph{Recall:} say the French for stage fright, then the French for line, then say \emph{interpréter} again
\end{itemize}

\begin{tcolorbox}[breakable,colback=teal!4,colframe=teal!35!black,title={Before you move on}]
What does interpréter mean? (\textquotedblleft{}To perform, to play (a role); to interpret\textquotedblright{}.) Say it once more.
\end{tcolorbox}
\section[le comédien]{le comédien — actor (especially on stage)}

\label{lesson:FR-C515-comedien}



\begin{quote}
Before the new one: say the French for line, then the French for to perform.
\end{quote}

\subsection*{You'll want to know: le comédien}
\textbf{le comédien} — \textquotedblleft{}actor (especially on stage)\textquotedblright{}.

\begin{quote}
Ce comédien joue au théâtre et au cinéma.
\end{quote}

\begin{quote}
\emph{This actor works in the theatre and in film.}
\end{quote}

A woman is \emph{la comédienne}. Not only comic actors: any actor.

\begin{cousinweb}
\emph{le comédien} is built on \emph{comédie} (EtymDB), from Latin \emph{cōmoedia}, a borrowing from Greek \emph{kōmōidía} \textquotedblleft{}comedy\textquotedblright{}.
\end{cousinweb}

\begin{grammarlens}[title={What you've built}]
One more word for telling, explaining and describing at B1.
\end{grammarlens}

\subsection*{Your turn}
\begin{itemize}
\raggedright
  \item \emph{Say it:} \emph{le comédien}
  \item \emph{Say it:} \emph{le comédien}, once more
  \item \emph{Say it:} a sentence of your own with \emph{le comédien}
  \item \emph{Recall:} say the French for line, then the French for to perform, then say \emph{le comédien} again
\end{itemize}

\begin{tcolorbox}[breakable,colback=teal!4,colframe=teal!35!black,title={Before you move on}]
What does le comédien mean? (\textquotedblleft{}Actor (especially on stage)\textquotedblright{}.) Say it once more.
\end{tcolorbox}
\section[les applaudissements]{les applaudissements — applause}

\label{lesson:FR-C515-applaudissements}



\begin{quote}
Before the new one: say the French for to perform, then the French for actor.
\end{quote}

\subsection*{You'll want to know: les applaudissements}
\textbf{les applaudissements} — \textquotedblleft{}applause\textquotedblright{}.

\begin{quote}
À la fin, il y a eu de longs applaudissements.
\end{quote}

\begin{quote}
\emph{At the end there was long applause.}
\end{quote}

Usually plural. The verb is \emph{applaudir} \textquotedblleft{}to clap\textquotedblright{}.

\begin{cousinweb}
\emph{les applaudissements} is built on \emph{applaudir} \textquotedblleft{}to clap\textquotedblright{}, a learned borrowing from Latin \emph{applaudere} \textquotedblleft{}to clap at\textquotedblright{}, built on \emph{plaudere} \textquotedblleft{}to clap\textquotedblright{} (Lewis \& Short).
\end{cousinweb}

\begin{grammarlens}[title={What you've built}]
One more word for telling, explaining and describing at B1.
\end{grammarlens}

\subsection*{Your turn}
\begin{itemize}
\raggedright
  \item \emph{Say it:} \emph{les applaudissements}
  \item \emph{Say it:} \emph{les applaudissements}, once more
  \item \emph{Say it:} a sentence of your own with \emph{les applaudissements}
  \item \emph{Recall:} say the French for to perform, then the French for actor, then say \emph{les applaudissements} again
\end{itemize}

\begin{tcolorbox}[breakable,colback=teal!4,colframe=teal!35!black,title={Before you move on}]
What does les applaudissements mean? (\textquotedblleft{}Applause\textquotedblright{}.) Say it once more.
\end{tcolorbox}
